\chapter{Telemetry, Tracking and Command (TT\&C)}

\newcommand{\ttcfield}[1]{{\ttfamily\footnotesize\detokenize{#1}}}

\section{Subsystem Overview}

The Telemetry, Tracking and Command (TT\&C) subsystem is the operational link
between Lakitu and its ground crew. It periodically downlinks the current
spacecraft state, provides the GNSS information used to track and recover the
balloon payload, accepts a small set of uplink messages, and exposes radio
health to Fault Detection, Isolation and Recovery (FDIR). The implemented link
is bidirectional half-duplex LoRa using an RFM95W/SX1276-class modem at
868~MHz. The flight computer is an STM32L476RG and the ground radio is driven
from a PC through a CH347 USB-to-SPI adapter.

The complete end-to-end system is shown in Figure~\ref{fig:ttc-architecture}.
An important design boundary is that Flight Software (FSW) owns construction
of the telemetry packet, while TT\&C owns radio scheduling, delivery attempts,
acknowledgement handling and modem health. FDIR, rather than TT\&C, owns the
policy that decides when a radio failure is severe enough to request recovery.

\begin{figure}[H]
\centering
\resizebox{0.97\textwidth}{!}{%
\begin{tikzpicture}[
  box/.style={draw, rounded corners=2pt, minimum height=1.15cm,
    minimum width=3.15cm, align=center,
    fill=blue!6, font=\small},
  radio/.style={box, fill=orange!14}, ground/.style={box, fill=green!10},
  store/.style={box, fill=violet!10}, flow/.style={-{Latex[length=2.5mm]}, thick},
  data/.style={font=\scriptsize, align=center, fill=white, inner sep=1.5pt}
]
  \node[box] (sensors) {Sensors, payload\\and SCV health};
  \node[box, right=1.0cm of sensors] (fsw) {FSW packet builder\\92-byte protocol v8};
  \node[box, right=1.0cm of fsw] (ttc) {TT\&C service\\queue, retry, ACK,\\radio health};
  \node[radio, right=1.9cm of ttc] (flight) {Flight RFM95W\\SPI1, 868 MHz};
  \node[radio, below=1.75cm of flight] (groundradio) {Ground RFM95W\\CH347 USB--SPI};
  \node[ground, left=1.9cm of groundradio] (backend) {Python/FastAPI\\decoder and API};
  \node[ground, left=1.9cm of backend] (ui) {React dashboard\\human-readable data};
  \node[store, below=0.75cm of backend] (csv) {CSV logs\\telemetry and events};
  \draw[flow] (sensors) -- node[above,data]{1 Hz} (fsw);
  \draw[flow] (fsw) -- node[above,data]{frame} (ttc);
  \draw[flow] (ttc) -- node[above,data]{non-blocking\\SPI} (flight);
  \draw[flow,<->] (flight) --
    node[right=2pt,data]{telemetry / ACKs /\\commands} (groundradio);
  \draw[flow] (groundradio) -- node[above,data]{payload /\\RSSI / SNR} (backend);
  \draw[flow] (backend) -- node[above,data]{WebSocket /\\REST} (ui);
  \draw[flow] (backend) -- (csv);
\end{tikzpicture}}
\caption{End-to-end TT\&C architecture and ownership boundaries.}
\label{fig:ttc-architecture}
\end{figure}

\subsection{Responsibilities}

The implemented subsystem configures and verifies the modem, maintains
continuous receive outside transmissions, schedules a nominal frame every
20~s and on significant events, and transmits without blocking the spacecraft
superloop. It retains an unacknowledged frame and retries the exact same bytes,
sequence and CRC up to three total attempts. It also parses acknowledgements
and commands, rejects malformed or replayed traffic, exposes raw health to
FDIR, and executes asynchronous isolation and recovery actions.

On the ground, TT\&C receives and validates frames, restores engineering units,
records link metadata and events, acknowledges valid downlinks, and provides
live and historical data to the operator dashboard. Tracking is not a separate
RF ranging function: it comes from the GNSS latitude, longitude, altitude,
speed, course and time carried in every telemetry frame.
\section{Requirements and Traceability}

Table~\ref{tab:ttc-requirements} maps the allocated requirements to the
present implementation. ``Implemented'' means that behavior exists in the
repository; it does not replace the environmental, range or compliance test
named by the requirement.

\begin{longtable}{p{1.25cm}p{5.0cm}p{5.5cm}p{2.3cm}}
\caption{TT\&C requirement traceability.}\label{tab:ttc-requirements}\\
\toprule
\textbf{ID} & \textbf{Requirement intent} & \textbf{Implementation evidence} & \textbf{Verification state} \\
\midrule
\endfirsthead
\toprule
\textbf{ID} & \textbf{Requirement intent} & \textbf{Implementation evidence} & \textbf{Verification state} \\
\midrule
\endhead
\bottomrule
\endfoot
MR-004 & Telemetry to 40~km altitude / 125~km worst-case slant range & Section~\ref{sec:ttc-link-budget} predicts 13.5~dB downlink margin with the workbook antenna assumptions. & Analysis passes; range test open \\
FR-022 & Approximately one packet every 20~s & \ttcfield{TTC_TELEMETRY_INTERVAL_MS = 20000}; event packets are also supported. & Timing implemented; duty-cycle issue open \\
FR-023 & Standard packet with state, position, IMU, payload, battery and integrity check & Protocol v8 carries all named data plus health; CRC-16 protects bytes 0--89. & Implemented and decoder-tested \\
FR-024 & European 868~MHz ISM-band operation & Flight and ground both select 868.000~MHz. & Implemented; hardware verification open \\
FR-025 & Receive, decode and display human-readable telemetry & RFM95W receiver, decoder/store/API, WebSocket and React views are present. & Implemented; hardware demonstration open \\
FR-026 & SF9, 125~kHz bandwidth, coding rate 4/5 & Flight and ground configure explicit-header LoRa, SF9, BW125 and CR~4/5. The stated $-137$~dBm sensitivity is not valid for SF9. & Registers implemented; sensitivity corrected \\
PR-004 & At least 10~dB margin at 125~km & The implemented powers predict 13.5~dB downlink and 9.5~dB uplink margin. & Downlink passes; uplink open \\
PR-007 & Margin $>10$~dB, raw SNR $>0$~dB and processed SNR $>10.5$~dB & The same model predicts raw SNR of +0.5~dB downlink and $-3.5$~dB uplink. & Downlink marginal; uplink fails \\
IF-005 & OBC-to-radio SPI no faster than 10~MHz & SPI1 is 625~kbit/s with interrupt-driven transfers. & Implemented \\
IF-007 & FSW supplies phase and system state in a defined format & \ttcfield{FSW_BuildTelemetryPacket()} consumes datapool, SCV and TTC health. & Implemented \\
CR-005 & Power and channel-use constraints & Flight PA is +14~dBm; the 20~s/513~ms schedule does not meet the project's stated 1\% arithmetic. & Duty-cycle closure open \\
SR-004 & No interference with balloon flight-control systems & No interference evidence is stored in the repository. & Operator test open \\
\end{longtable}

\section{Physical and Radio Interface}

The RFM95W is connected to STM32 SPI1 in master mode. Transfers are full
duplex, 8-bit and interrupt driven. Chip select and reset are ordinary GPIOs;
the driver polls the modem IRQ register, so the implemented path does not need
a DIO interrupt line.

\begin{table}[H]
\centering
\begin{tabular}{llll}
\toprule
\textbf{Function} & \textbf{STM32 pin} & \textbf{Board label} & \textbf{RFM95W signal} \\
\midrule
SPI clock & PA5 & D13 & SCK \\
SPI MISO & PA6 & D12 & MISO \\
SPI MOSI & PA7 & D11 & MOSI \\
Chip select & PB6 & D10 & NSS / CS \\
Reset & PC7 & D9 & RESET \\
\bottomrule
\end{tabular}
\caption{Flight radio electrical interface.}
\label{tab:ttc-pins}
\end{table}

The CubeMX configuration and hand-maintained initialization both set SPI1 to
625~kbit/s. The driver verifies the RFM95W version register (expected
\texttt{0x12}) and reads back frequency, modem and sync-word registers. A
transport failure produces \texttt{INIT\_FAIL}; a readable but unexpected
value produces \texttt{CONFIG\_FAIL}.

\begin{table}[H]
\centering
\begin{tabular}{p{4.2cm}p{4.0cm}p{5.5cm}}
\toprule
\textbf{Parameter} & \textbf{Flight value} & \textbf{Ground implementation} \\
\midrule
Carrier frequency & 868.000~MHz (FRF \texttt{D9 00 00}) & Backend default 868.000~MHz \\
Spreading factor / bandwidth & SF9 / 125~kHz & SF9 / 125~kHz \\
Coding rate / header & 4/5 / explicit & 4/5 / explicit \\
Radio payload CRC / preamble & Enabled / 8 symbols & Enabled / 8 symbols \\
Sync word & \texttt{0x12} & \texttt{0x12} \\
Transmit power & +14~dBm (PA boost) & +10~dBm default \\
Receive mode & Continuous between TX & Continuous except during TX \\
Telemetry payload & 92 bytes & Accepts exactly 92 bytes \\
\bottomrule
\end{tabular}
\caption{LoRa configuration derived from flight and ground source code.}
\label{tab:ttc-radio-config}
\end{table}

Flight and ground now use the same 868.000~MHz center frequency. An explicit
pre-flight register/readback check remains recommended to verify the deployed
configuration before operation.
\section{Telemetry Generation and Scheduling}

The superloop calls \ttcfield{TTC_Service()} on every iteration because TT\&C
and the LoRa driver are cooperative, non-blocking state machines. They advance
at most a small SPI or state step per call, allowing FDIR and the independent
watchdog to continue while a radio operation is in progress.

When \ttcfield{TTC_TelemetryDue()} is true, FSW builds a packet from the 1~Hz
sensor datapool, persistent Spacecraft Vector (SCV), and volatile TT\&C
health/command snapshots. A frame is due before the first transmission, after
20~s from the last successful TX, or when an immediate request is latched.
Flight-phase changes and received commands set this latch. If a previous frame
still awaits its ACK, no new frame is built: the pending sequence is completed
or abandoned first.

Current sensor validity is separate from the numeric value. If a source becomes
invalid, its last known value is retained but its validity bit is cleared.
Before the first valid observation, signed fields use the minimum integer and
unsigned fields use their maximum. Ground must therefore interpret validity,
the FDIR fault mask and the equipment-enabled mask together.

\section{Protocol-v8 Telemetry Frame}

The downlink is a packed, little-endian, fixed-width 92-byte binary frame.
Engineering values are scaled and clamped to integer units; ground reverses the
scaling. Figure~\ref{fig:ttc-packet-groups} shows the allocation and
Table~\ref{tab:ttc-wire-layout} gives the complete definition.

\begin{figure}[H]
\centering
\begin{tikzpicture}[x=0.15cm,y=1cm,font=\scriptsize]
  \draw[fill=blue!14] (0,0) rectangle (26,1) node[midway,align=center]{Envelope, time, phase,\\validity and FDIR\\bytes 0--25};
  \draw[fill=green!16] (26,0) rectangle (50,1) node[midway,align=center]{GNSS navigation\\bytes 26--49};
  \draw[fill=orange!18] (50,0) rectangle (62,1) node[midway,align=center]{IMU\\50--61};
  \draw[fill=cyan!16] (62,0) rectangle (74,1) node[midway,align=center]{Baro + batt.\\62--73};
  \draw[fill=violet!16] (74,0) rectangle (81,1) node[midway,align=center]{Coral\\74--80};
  \draw[fill=red!13] (81,0) rectangle (90,1) node[midway,align=center]{Radio /\\uplink\\81--89};
  \draw[fill=gray!25] (90,0) rectangle (92,1);
  \draw[-{Latex},thin] (91,0) -- (91,-0.35) -- (84,-0.35) node[left]{CRC-16, bytes 90--91};
  \draw[decorate,decoration={brace,amplitude=4pt},thick]
       (0,1.15) -- (92,1.15) node[midway,yshift=9pt]{92 bytes total};
\end{tikzpicture}
\caption{Protocol-v8 byte allocation by information group.}
\label{fig:ttc-packet-groups}
\end{figure}

\begin{longtable}{r r p{4.2cm} p{7.0cm}}
\caption{Complete protocol-v8 telemetry wire layout.}\label{tab:ttc-wire-layout}\\
\toprule
\textbf{Offset} & \textbf{Bytes} & \textbf{Wire field} & \textbf{Meaning and scale} \\
\midrule
\endfirsthead
\toprule
\textbf{Offset} & \textbf{Bytes} & \textbf{Wire field} & \textbf{Meaning and scale} \\
\midrule
\endhead
\bottomrule
\endfoot
0 & 1 & \ttcfield{packet_type} & \texttt{0x01} identifies telemetry \\
1 & 1 & \ttcfield{protocol_version} & \texttt{0x08}; ground accepts v8 only \\
2 & 2 & \ttcfield{sequence_number} & Unsigned sequence; unchanged on retry \\
4 & 4 & \ttcfield{tx_uptime_s} & OBC uptime in seconds \\
8 & 4 & \ttcfield{mission_elapsed_s} & Persistent mission time in seconds \\
12 & 2 & \ttcfield{boot_count_sat} & Boot count, saturated at 65,535 \\
14 & 1 & \ttcfield{flight_phase} & Six-state flight-phase enum \\
15 & 1 & \ttcfield{reset_reason} & Unknown, power-on, watchdog or software reset \\
16 & 1 & \ttcfield{validity_flags} & Bits 0--4: GNSS, IMU, baro, battery, Coral \\
17 & 1 & \ttcfield{i2c_bus_state} & 0 idle, 1 restart, 2 hold \\
18 & 2 & \ttcfield{equipment_enabled} & Policy mask for enabled equipment \\
20 & 2 & \ttcfield{equipment_faults} & Debounced FDIR fault mask \\
22 & 2 & \ttcfield{sample_age_ms_sat} & Packet tick minus datapool tick; 65,535 unknown/overflow \\
24 & 1 & \ttcfield{watchdog_reset_count} & Persisted watchdog-reset count \\
25 & 1 & \ttcfield{sd_fault_count} & Persisted SD fault count \\
26 & 4 & \ttcfield{latitude_e7} & Latitude, degrees $\times 10^7$ \\
30 & 4 & \ttcfield{longitude_e7} & Longitude, degrees $\times 10^7$ \\
34 & 4 & \ttcfield{gnss_altitude_dm} & GNSS altitude in decimetres \\
38 & 2 & \ttcfield{vertical_speed_cms} & cm/s, positive upward \\
40 & 2 & \ttcfield{ground_speed_cms} & Horizontal speed in cm/s \\
42 & 2 & \ttcfield{course_cdeg} & Course in centidegrees, 0--35,999 \\
44 & 4 & \ttcfield{gnss_utc_sod} & UTC seconds of day \\
48 & 1 & \ttcfield{gnss_satellites} & Satellites used in solution \\
49 & 1 & \ttcfield{gnss_fix_type} & GNSS fix-type enum \\50 & 2 & \ttcfield{accel_x_mg} & X acceleration in milli-g \\
52 & 2 & \ttcfield{accel_y_mg} & Y acceleration in milli-g \\
54 & 2 & \ttcfield{accel_z_mg} & Z acceleration in milli-g \\
56 & 2 & \ttcfield{gyro_x_ddeg_s} & X angular rate in 0.1 degrees/s \\
58 & 2 & \ttcfield{gyro_y_ddeg_s} & Y angular rate in 0.1 degrees/s \\
60 & 2 & \ttcfield{gyro_z_ddeg_s} & Z angular rate in 0.1 degrees/s \\
62 & 4 & \ttcfield{baro_pressure_pa} & Barometric pressure in pascals \\
66 & 4 & \ttcfield{baro_altitude_dm} & Ground-relative altitude in decimetres \\
70 & 2 & \ttcfield{baro_temperature_cdeg} & Temperature in 0.01 degrees C \\
72 & 2 & \ttcfield{battery_mv} & Battery voltage in millivolts \\
74 & 2 & \ttcfield{coral_sequence_low} & Low 16 bits of payload result sequence \\
76 & 2 & \ttcfield{coral_fraction_q16} & Cloud fraction, 0--65,535 maps to 0--100\% \\
78 & 1 & \ttcfield{coral_status} & Coral status bit mask \\
79 & 2 & \ttcfield{coral_age_s_sat} & Coral result age in seconds; saturated \\
81 & 1 & \ttcfield{lora_last_event} & Latest modem event enum \\
82 & 1 & \shortstack[l]{\ttcfield{lora_consecutive_}\\\ttcfield{failures}} & Consecutive radio failures \\
83 & 1 & \ttcfield{lora_recovery_count} & Recovery attempts since boot \\
84 & 1 & \ttcfield{lora_rx_state} & Bits 0--2 RX status; bit 3 continuous RX active \\
85 & 1 & \shortstack[l]{\ttcfield{lora_tx_fault_}\\\ttcfield{count_sat}} & Lifetime failed TX attempts, saturated \\
86 & 1 & \shortstack[l]{\ttcfield{lora_ack_timeout_}\\\ttcfield{count_sat}} & Exhausted ACK budgets, saturated \\
87 & 2 & \ttcfield{last_command_id} & Last command ID observed by flight \\
89 & 1 & \ttcfield{uplink_state} & Bits 0--2 command status; 3--5 ACK status \\
90 & 2 & \ttcfield{crc16} & CRC-16/CCITT-FALSE over bytes 0--89 \\
\end{longtable}

The CRC uses polynomial \texttt{0x1021} and initial value \texttt{0xFFFF}.
Ground first rejects any payload not exactly 92 bytes, then checks packet type,
protocol version and application CRC. The RFM95W packet CRC is an additional
physical-link check. Sequence numbers identify gaps, duplicates and a new
session after reboot.

Protocol v8 reduced the former 155-byte raw-structure frame to 92 bytes, a 41\%
payload reduction. Estimated airtime fell from about 800~ms to 513~ms, about
36\% per attempt. Data without an onboard producer, such as GNSS dilution of
precision and MCU/IMU temperature, deliberately receive no wire bytes.

\section{Reliable Downlink and Uplink Protocol}

LoRa supplies the radio packet, while reliability is implemented at the
application layer. Figure~\ref{fig:ttc-reliability} summarizes a transaction.

\begin{figure}[H]
\centering
\resizebox{0.94\textwidth}{!}{%
\begin{tikzpicture}[
  flowbox/.style={draw, rounded corners, minimum width=4.3cm,
    minimum height=0.9cm, align=center, fill=blue!7},
  gstep/.style={flowbox, fill=green!10},
  decision/.style={diamond, draw, aspect=2.2, align=center,
    fill=orange!12, inner sep=1pt},
  flow/.style={-{Latex[length=2.5mm]}, thick},
  note/.style={font=\scriptsize, align=center, fill=white, inner sep=1.5pt}
]
  \node[flowbox] (build) {Flight: retain sequence $N$};
  \node[flowbox, below=0.85cm of build] (tx)
    {Flight: transmit exact 92-byte frame};
  \node[gstep, right=3.2cm of tx] (validate)
    {Ground: verify length, v8 and CRC};
  \node[gstep, below=0.85cm of validate] (ack)
    {Ground: log/display, send \texttt{ACK,N}};
  \node[decision, below=1.60cm of tx] (got) {Matching ACK\\within 5 s?};
  \node[decision, below=1.25cm of got] (tries) {Three attempts\\used?};
  \node[flowbox, below right=0.45cm and 1.8cm of got] (done)
    {Release frame\\resume schedule};
  \node[flowbox, right=1.8cm of tries] (abandon)
    {Count ACK timeout\\release frame};
  \draw[flow] (build) -- (tx);
  \draw[flow] (tx) -- node[above,note]{downlink $N$} (validate);
  \draw[flow] (validate) -- (ack);
  \draw[flow] (ack.west) -- ++(-0.65,0) -|
    node[pos=0.55,above,note]{uplink ACK} ([xshift=-0.8cm]got.west)
    -- (got.west);
  \draw[flow] (got.east) -| node[pos=0.75,right,note]{yes} (done.north);
  \draw[flow] (got) -- node[right,note]{no} (tries);
  \draw[flow] (tries.east) -- node[above,note]{yes} (abandon.west);
  \draw[flow] (tries.south) -- ++(0,-0.85) -|
    node[pos=0.25,below,note]{no: retry same bytes} ([xshift=-1.1cm]tx.west)
    -- (tx.west);
\end{tikzpicture}}
\caption{Application-level acknowledgement and retry sequence.}
\label{fig:ttc-reliability}
\end{figure}

After \textit{TxDone}, flight returns to continuous RX and waits 5~s for
\texttt{ACK,<sequence>}. A match releases the frame; the last accepted ACK is
classified as duplicate; another valid sequence is unexpected. ACKs require no
ACK-of-ACK. Ground sends the ACK only after full application validation and
acknowledges a duplicated sequence again. A duplicate normally means downlink
succeeded but the earlier ground ACK did not reach flight in time.
\subsection{Commands and Replay Protection}

Flight accepts ASCII uplinks up to 64 bytes. The implemented classes are:

\begin{table}[H]
\centering
\begin{tabular}{p{4.5cm}p{4.7cm}p{5.0cm}}
\toprule
\textbf{Wire form} & \textbf{Meaning} & \textbf{Flight response} \\
\midrule
\texttt{ACK,<sequence>} & Acknowledge telemetry & Update independent ACK-status latch \\
\ttcfield{CMD,<id>,REQ_TELEMETRY} & Request out-of-cycle telemetry & Record ID/status and schedule immediate telemetry \\
\ttcfield{CMD,<id>,<other verb>} & Valid but unsupported command & Latch \texttt{UNSUPPORTED} and schedule telemetry \\
Malformed \texttt{CMD} or \texttt{ACK} & Invalid format & Latch \texttt{INVALID\_FORMAT} and schedule telemetry \\
\bottomrule
\end{tabular}
\caption{Implemented uplink message classes.}
\label{tab:ttc-uplink}
\end{table}

Command IDs are non-zero 16-bit values. Flight maintains a 16-command sliding
replay window with wrap-safe arithmetic: a command is accepted once and a
repeated ID is reported as duplicate. Ground serializes command transactions,
assigns a stable ID, and by default retries three times with 5~s TX and 5~s
flight-response timeouts. Confirmation is the command ID and
\texttt{ACCEPTED}/\texttt{DUPLICATE} status echoed in valid telemetry.
Command and telemetry-ACK status are independent latches, so an ACK cannot erase
a command confirmation. \texttt{REQ\_TELEMETRY} is currently the only
operational verb.

\section{Flight Software State Machines}

The TT\&C layer and LoRa driver are asynchronous. Initialization, send,
receive-start and recovery calls queue work, and repeated service calls advance
it. This prevents radio waits from starving the watchdog or other subsystems.

\begin{table}[H]
\centering
\begin{tabular}{p{3.7cm}p{10.5cm}}
\toprule
\textbf{State family} & \textbf{Purpose} \\
\midrule
Startup initialization / RX & Reset/configure, verify registers, enter continuous RX \\
Idle & Poll uplink, reconcile faults, start telemetry/retries and FDIR actions \\
TX wait / RX restart & Await \textit{TxDone}/timeout, record result, restore RX \\
Action isolation & Quiesce SPI and hold modem in reset \\
Action initialization / RX & Full recovery or return-to-service, then restore RX \\
\bottomrule
\end{tabular}
\caption{High-level TT\&C state-machine responsibilities.}
\end{table}

Driver transfers use \ttcfield{HAL_SPI_TransmitReceive_IT()}; modem IRQ flags
are polled for \textit{RxDone}, CRC error and \textit{TxDone}. SPI timeouts
trigger a bounded interrupt abort, and the driver can reinitialize SPI if it
does not quiesce. Oversized observations are consumed and returned as length
errors so stale FIFO data cannot poison the next packet.

\section{FDIR and Radio Health}

TT\&C records facts; FDIR applies policy. The raw snapshot contains readiness,
continuous-RX state, consecutive/lifetime failures, successful TX/RX times,
recovery activity, isolation and exhausted-ACK count. FDIR reconstructs a
pass/fail history from the counters.

\begin{table}[H]
\centering
\begin{tabular}{p{4.3cm}p{5.2cm}p{4.7cm}}
\toprule
\textbf{Condition / action} & \textbf{Implemented policy} & \textbf{Ownership} \\
\midrule
Recovery request & 5 consecutive failures, or 12 failures in last 20 attempts & FDIR decides; TT\&C executes \\
Recovery cooldown & At least 10~s between requests & FDIR \\
SCV LoRa fault & Set at 9 failures in last 10; clears as window improves & FDIR \\
Full recovery & Reset, configure, read back, restart RX & TT\&C / driver \\
RX restart & Re-enter RX without full reset if otherwise ready & TT\&C / driver \\
Isolation & Preempt activity and hold reset & API exists; automatic give-up policy open \\
Return to service & Reinitialize isolated/faulted modem & TT\&C / driver \\
\bottomrule
\end{tabular}
\caption{Implemented LoRa FDIR policy and ownership.}
\label{tab:ttc-fdir}
\end{table}

FDIR sets \ttcfield{EQUIPMENT_LORA} in a shared reinit bitmask.
\ttcfield{TTC_Service()} consumes it and advances the action later. Only one
action may be pending, except isolation can preempt. Requests are idempotent.
Telemetry reports the latest event, failures, recoveries, RX state, lifetime TX
faults and ACK timeouts. Ground combines this snapshot with end-to-end evidence
such as a received frame, RSSI/SNR, duplicates and command outcomes.

\section{Ground Segment}

The ground segment has three layers. Python controls an RFM95W through CH347,
reads payload/RSSI/SNR/CRC, and restores continuous RX after transmissions.
FastAPI validates and decodes v8, tracks lost/duplicate sequences, keeps bounded
history, writes telemetry and event CSVs, exposes REST, and broadcasts over
WebSocket. The React dashboard provides Overview, Health, Position/Flight
Phase, Link Analytics, Commands and Raw Telemetry views.

Invalid and physical-CRC observations remain available for diagnosis but never
replace the latest valid spacecraft state. Main interfaces are
\texttt{/api/status}, \texttt{/api/latest}, \texttt{/api/history},
\texttt{/api/stats}, \texttt{/api/uplink-log}, \texttt{/api/send-ascii} and
\texttt{/ws/telemetry}. The backend can run without radio for UI/API work.
User verbs are limited to 48 ASCII characters so the envelope fits flight RX.
\section{Link Budget, Airtime and Channel Use}
\label{sec:ttc-link-budget}

The link-budget workbook uses 35~km altitude and 120~km horizontal drift, giving
125~km slant range. MR-004 instead states 40~km altitude with the same drift,
which would be 126.5~km; this adds only 0.1~dB path loss but should be reconciled.
Free-space path loss is
\begin{equation}
L_{\mathrm{FS}} = 32.44 + 20\log_{10}(d_{\mathrm{km}})
                 + 20\log_{10}(f_{\mathrm{MHz}}),
\end{equation}
which gives 133.15~dB at 125~km and 868~MHz. The workbook assumes +20~dBm
flight transmit power, although the implemented PA configuration is +14~dBm.
It also contains a $-130$~dBm sensitivity cell while its note and FR-026 state
$-137$~dBm. For SF9/BW125, the workbook noise model and RFM95W demodulator
threshold give
\begin{align}
N &= -174 + 10\log_{10}(125000) + 6 = -117.0~\mathrm{dBm},\\
S_{\min} &= N + (-12.5) = -129.5~\mathrm{dBm}\approx-130~\mathrm{dBm}.
\end{align}
The RFM95W data sheet associates approximately $-136$~dBm with the slower
SF12/BW125 setting, not SF9. Table~\ref{tab:ttc-link-budget} therefore uses
$-130$~dBm and retains the workbook's antenna, cable, atmospheric and
polarization assumptions. It is an analytical estimate, not RF qualification.

\begin{table}[H]
\centering
\begin{tabular}{lrr}
\toprule
\textbf{Term} & \textbf{Downlink} & \textbf{Uplink} \\
\midrule
Transmit power & +14.00~dBm & +10.00~dBm \\
Transmit antenna gain & +2.15~dBi & +5.00~dBi \\
Transmit cable loss & $-0.50$~dB & $-0.50$~dB \\
Free-space path loss & $-133.15$~dB & $-133.15$~dB \\
Atmospheric loss & $-0.50$~dB & $-0.50$~dB \\
Polarization loss & $-3.00$~dB & $-3.00$~dB \\
Receive antenna gain & +5.00~dBi & +2.15~dBi \\
Receive cable loss & $-0.50$~dB & $-0.50$~dB \\
\textbf{Predicted receive power} & \textbf{$-116.50$~dBm} & \textbf{$-120.50$~dBm} \\
SF9 sensitivity & $-130.00$~dBm & $-130.00$~dBm \\
\textbf{Link margin} & \textbf{13.50~dB} & \textbf{9.50~dB} \\
Noise floor & $-117.01$~dBm & $-117.01$~dBm \\
\textbf{Raw SNR} & \textbf{+0.51~dB} & \textbf{$-3.49$~dB} \\
\bottomrule
\end{tabular}
\caption{Recomputed 125~km link budget using implemented transmit powers and workbook RF assumptions.}
\label{tab:ttc-link-budget}
\end{table}

Using the implemented +14~dBm flight power reduces the workbook's downlink
receive power, link margin and raw SNR by 6~dB. The resulting downlink meets the
10~dB margin target by 3.5~dB and its modeled raw SNR is only 0.5~dB above
PR-007. The +10~dBm ground transmitter gives 9.5~dB uplink margin and
$-3.5$~dB raw SNR, so the bidirectional command/ACK path is not demonstrated at
the worst-case range. The requirements file's 26.5~dB claim is not reproduced.
Acceptance requires measured sensitivity, characterized installed antennas and
cables, and representative bidirectional range testing. The actual ground
antenna gain must also be confirmed: the workbook note says 3--4~dBi while its
calculation uses 5~dBi.

At SF9, BW125, CR~4/5, explicit header, CRC and eight-symbol preamble, symbol
time is 4.096~ms and a 92-byte frame uses 113 payload symbols, giving
513.024~ms airtime. The current packet therefore carries $92\times8/20=36.8$
payload bit/s; this is 2.09\% of the nominal 1758~bit/s LoRa data rate, but it
does not represent transmitter channel occupancy. At a strict 20~s period:
\begin{equation}
\frac{0.513}{20}\times100\% \approx 2.57\%.
\end{equation}
This conflicts with the 1\% channel-use requirement in FR-022/CR-005 before
event packets or retries. At the same modulation, one 513~ms transmission needs
at least a 51.3~s interval to average 1\%. The project must close this by a
permitted channel/access regime, schedule/modulation change, or another
reviewed compliance basis. Retries and ground ACKs belong in the final analysis.
The workbook's data-budget sheet still uses the obsolete 76-byte packet and
must be updated to 92 bytes before it is used as verification evidence.

\section{Verification Evidence}

Repository tests provide the following host-side evidence:
\begin{itemize}[leftmargin=*]
  \item v8 decoder tests cover size, scaling, sentinels, CRC and legacy rejection;
  \item TT\&C harnesses cover startup recovery, RX faults, idempotency,
        isolation, retries, command/ACK latches and the FDIR request bit;
  \item driver harnesses cover oversized RX, 255-byte FIFO transfer, abort
        rejection/timeout recovery and abort callbacks;
  \item source-contract tests enforce interrupt-driven SPI, non-blocking design,
        FDIR policy ownership and queued recovery interfaces.
\end{itemize}

Remaining acceptance work is a two-radio register dump; end-to-end frame,
ACK and command demonstration; five-minute schedule measurement; injected CRC,
SPI, ACK-loss and reset faults; sensitivity and antenna characterization;
range test; actual channel-occupancy log including retries; and an interference
check coordinated with the balloon operator.

\section{Design Decisions and Known Limitations}

The central decisions are a portable fixed-integer frame, continuous receive,
reliability and replay protection above LoRa, non-blocking state machines,
independent command/ACK latches, and FDIR ownership of recovery policy. These
improve watchdog safety, diagnostics and operational observability.

Open limitations are:
\begin{itemize}[leftmargin=*,nosep]
  \item flight and ground center frequencies differ by 100~kHz;
  \item 20~s scheduling and 513~ms airtime conflict with the stated 1\% target;
  \item range margin depends on unverified sensitivity and unspecified antennas;
  \item \texttt{REQ\_TELEMETRY} is the only implemented command verb;
  \item automatic give-up escalation to long-term isolation remains open;
  \item operator/deployment procedures need a dedicated ground-station runbook;
  \item no artifact yet proves 125~km range, regulatory closure or
        non-interference.
\end{itemize}

Closing these items is primarily integration and verification work; the packet,
reliable transactions, decoder and core radio/FDIR state machines exist.
